% =====================================================================
\chapter{Discusión}\label{cap:discusion}
% Plantilla ITBA: relacionar resultados con teoría y estado del arte,
% contrastar con la hipótesis, implicancias y limitaciones.
% Esta version cubre NOR, SLR y OF. El resto de Resultados
% (consumo/peso, neurogenesis, ML) queda pendiente, ver nota al final.
% =====================================================================

\section{Síntesis de los hallazgos}\label{sec:disc-sintesis}

La pregunta que guía este trabajo (\S\ref{sec:planteamiento}) indaga si el consumo temprano de stevia reproduce los efectos adversos documentados para la sacarosa sobre la memoria de reconocimiento y la separación de patrones espaciales, y en qué medida ese efecto depende de la ventana etaria de exposición. Los resultados del NOR y la SLR ofrecen una respuesta parcial y asimétrica entre pruebas. La sacarosa, empleada como control positivo (\S\ref{sec:bases}), reprodujo el fenotipo adverso esperado: redujo la tasa de exploración del objeto novedoso de forma consistente en las tres demoras del NOR cuando la exposición fue juvenil, y redujo el índice de discriminación (DI) de la SLR en ambas configuraciones geométricas. La stevia, en cambio, no reprodujo ese patrón en ninguna de las dos pruebas, lo que es coherente con la primera mitad de la hipótesis (\S\ref{sec:hipotesis}): que la stevia no induciría los efectos adversos de la sacarosa. Sin embargo, los datos no permiten sostener la segunda mitad de la hipótesis, que postula una condición equivalente o superior a la del control. En la SLR, los grupos stevia no se diferenciaron del control pero tampoco lograron discriminar la posición novedosa por encima del azar, y en el NOR mostraron efectos puntuales e inconsistentes entre demoras. La hipótesis, por tanto, no queda refutada y la evidencia se inclina hacia su lado favorable: la stevia no reprodujo los efectos de la sacarosa, no se diferenció del control en el DI de la SLR ni de forma replicable en el NOR, y mostró en el \textit{open field} un perfil de mayor exploración de las zonas centrales. El perfil que emerge es el de un compuesto que, a diferencia de la sacarosa, no perturba de forma consistente la memoria evaluada, y que en las condiciones de este trabajo resulta más cercano a no perjudicial que a perjudicial. Esa lectura se acompaña de dos reservas: la equivalencia con el control no puede afirmarse con la potencia disponible (en la SLR los grupos stevia no discriminaron por encima del azar), y la ausencia de diferencia no equivale a una inocuidad demostrada. Los resultados del \textit{open field}, que aportan un perfil conductual adicional, se discuten en la sección~\ref{sec:disc-of}.

\section{Validez de las pruebas}\label{sec:disc-validez}

La interpretación de estos resultados depende, en primer lugar, de que las pruebas conductuales hayan funcionado como se espera. En este sentido, los grupos control discriminaron el objeto novedoso por encima del nivel de azar en las tres demoras del NOR, en ambas ventanas etarias (\S\ref{sec:res-nor-validacion}), y discriminaron la posición novedosa por encima del azar en ambas configuraciones de la SLR, también en ambas edades (\S\ref{sec:res-slr-validacion}). Esto confirma que la batería detecta memoria de reconocimiento a las demoras evaluadas \cite{leger2013nor} y separación de patrones espaciales bajo las dos demandas de carga cognitiva operacionalizadas \cite{reichelt2021slr}, y que la ausencia de maduración completa del circuito hipocampal en PD75 no impide, por sí sola, un desempeño por encima del azar.

Esta validación, sin embargo, es parcial. La exploración total del día 1 de la SLR, pensada como control de codificación, no fue equivalente entre grupos (\S\ref{sec:res-slr-exploracion}), y la exploración total durante la adquisición del NOR (T1) todavía no está analizada. En cambio, en el \textit{open field} la distancia total recorrida no difirió entre grupos (\S\ref{sec:res-of-locomocion}), lo que no sugiere un cambio de la actividad locomotora general, aunque el menor tamaño de los grupos expuestos a stevia limita esa conclusión. Mientras el control de T1 siga pendiente y la exploración del día 1 de la SLR no sea equivalente, no puede descartarse por completo que una diferencia de codificación o de exploración, y no específicamente de memoria, contribuya a los efectos observados en T2--T4 del NOR o en el día 2 de la SLR; los hallazgos de las secciones siguientes deben leerse con esa reserva.

\section{Sacarosa: el control positivo}\label{sec:disc-sacarosa}

En el NOR, la exposición juvenil a sacarosa redujo la tasa de exploración del objeto novedoso respecto del control juvenil en las tres demoras evaluadas, mientras que la sacarosa adulta no difirió de su control en ninguna de ellas; la interacción edad $\times$ tratamiento fue significativa en T2, T3 y T4 (\S\ref{sec:res-nor-sacarosa}). Esta disociación por edad sí puede afirmarse con el respaldo del test estadístico correspondiente, y reproduce de forma cercana el hallazgo de Coirini et~al., quienes documentaron que la exposición juvenil a sacarosa compromete la memoria a largo plazo con mayor severidad que la exposición adulta \cite{coirini2022sucrose}. También concuerda con trabajos en los que el acceso adolescente a sacarosa o a jarabe de maíz de alta fructosa se asoció con alteraciones de la memoria hipocampal \cite{hsu2015sucrose,reichelt2015adolescent}. Difiere, en cambio, de lo descrito por Noble et~al., quienes con consumo temprano de azúcar encontraron un deterioro de la memoria contextual dependiente del hipocampo con el reconocimiento simple de objetos preservado \cite{noble2019early}; en este trabajo, la sacarosa juvenil alteró el NOR incluso a las 2\,h. Las diferencias de protocolo entre estudios (tarea, demoras, tipo y concentración de la solución y tiempo transcurrido desde el retiro) pueden contribuir a esa discrepancia, pero esta tesis no permite distinguir cuál. La consistencia del efecto a través de demoras que reclutan de forma creciente al hipocampo \cite{leger2013nor} sugiere que el déficit juvenil no se restringe a una fase particular de la consolidación, sino que compromete el procesamiento de la novedad de forma más general en esa ventana etaria, en línea con la lógica de ventanas críticas que fundamenta el diseño (\S\ref{sec:ventanas}).

En la SLR, la sacarosa redujo el DI respecto del control con un efecto principal de tratamiento significativo en ambas configuraciones (\S\ref{sec:res-slr-sacarosa}), lo que extiende a la separación de patrones espaciales el fenotipo adverso ya descrito para la sacarosa sobre otras formas de memoria dependientes del hipocampo \cite{reichelt2016sucrose}. A diferencia del NOR, la interacción edad $\times$ tratamiento no alcanzó significancia en ninguna configuración, de modo que el patrón observado —sacarosa juvenil significativamente menor que el control en ambas configuraciones; sacarosa adulta significativamente menor solo en s-SLR— es descriptivo y no puede interpretarse como una disociación estadísticamente demostrada entre edades. Tampoco se evaluó formalmente la interacción configuración $\times$ tratamiento, por lo que no corresponde afirmar un déficit selectivo de separación de patrones en los adultos; lo que sí puede decirse es que el efecto de la sacarosa adulta alcanzó significancia únicamente en la configuración de mayor demanda de separación, un patrón consistente con —aunque no demostrativo de— un compromiso más marcado de los procesos dependientes de neurogénesis que del componente espacial general \cite{reichelt2016sucrose,perezgarcia2016malnutrition}. El patrón se asemeja a la disociación descrita en ratas tras una exposición corta a azúcar, que afectó la memoria de lugar y no el reconocimiento de objetos \cite{beilharz2014place} y, de forma selectiva, la memoria dependiente del hipocampo \cite{beilharz2016liquid}, aunque la exposición, la edad y las tareas de esos trabajos difieren de las de esta tesis.

Los mecanismos candidatos para este fenotipo —una reducción de la proliferación o la supervivencia de neuronas granulares nuevas, o una disrupción de la microbiota intestinal con efectos río abajo sobre la plasticidad hipocampal— tienen respaldo en la literatura \cite{maniam2016sugar,noble2021microbial}, a lo que se suman la neuroinflamación hipocampal \cite{hsu2015sucrose} y las alteraciones de la inmunorreactividad de parvalbúmina y doblecortina \cite{reichelt2015adolescent,xu2018sucrose}, aunque el cuadro no es uniforme: la exposición a sacarosa restringida a la adolescencia tardía llegó a aumentar la proliferación en la zona subgranular, pese a asociarse con neuronas inmaduras de morfología alterada \cite{sanchezhuerta2022sucrose}, lo que indica que el signo del efecto sobre la neurogénesis puede depender de la ventana exacta de exposición. En esta tesis esos mecanismos permanecen como hipótesis: la cuantificación de PCNA y DCX (\S\ref{sec:res-neurogenesis}) es la que podría determinar si el déficit conductual de la sacarosa se acompaña de una alteración en alguna etapa específica de la cascada neurogénica, y no antes.

\section{Stevia y contraste con la hipótesis}\label{sec:disc-stevia}

El resultado de la stevia en la SLR es el que exige mayor cautela interpretativa. En ninguna configuración ni edad el DI del grupo stevia difirió del control; pero tampoco en ninguna de ellas el grupo stevia discriminó la posición novedosa por encima del azar (\S\ref{sec:res-slr-stevia}), a diferencia de los controles, que sí lo hicieron en las cuatro combinaciones de edad y configuración. La ausencia de diferencia con el control no equivale, por tanto, a una preservación demostrada de la función: una interpretación igualmente compatible con los datos es que el grupo stevia ocupa una posición intermedia, no resuelta por el tamaño muestral disponible, entre el desempeño del control y el de la sacarosa. En la misma dirección, la variabilidad del DI de los grupos stevia fue considerablemente mayor que la de los controles —en el caso más marcado, la d-SLR juvenil, el error estándar del grupo stevia (0,100) duplicó al del control (0,050) (Tabla~\ref{tab:slr-di})—, lo que es compatible con una respuesta heterogénea entre animales más que con un efecto nulo homogéneo. Esta ambigüedad no puede resolverse con los datos actuales y debería tratarse como un resultado abierto, no como evidencia de inocuidad.

En el NOR, la stevia tampoco ofrece un patrón estable: en el grupo juvenil difirió del control únicamente a las 24\,h (T3), y en el adulto únicamente a las 2\,h (T2), sin que el efecto se sostuviera en las demoras restantes (\S\ref{sec:res-nor-stevia}). Esta inconsistencia admite al menos cuatro lecturas no excluyentes. Primera, que el efecto sea real pero dependiente de la fase de la memoria que predomina en cada demora —más dependiente de la corteza perirrinal a las 2\,h y más del hipocampo hacia las 24\,h y los 7 días \cite{leger2013nor}—, de modo que la stevia afectaría de forma distinta a cada sustrato según la edad de exposición. Segunda, que se trate de falsos positivos aislados: cada demora se analizó con un ANOVA de dos vías independiente seguido de comparaciones \textit{post hoc}, de modo que el número efectivo de contrastes realizados sobre el efecto de la stevia es mayor que uno, y la distinción entre un hallazgo significativo y uno no significativo en comparaciones de magnitud similar no es en sí misma necesariamente significativa \cite{gelman2006difference}. Tercera, que exista variabilidad entre las cohortes integradas para el análisis (\S\ref{sec:diseno}), de modo que un efecto puntual en una cohorte se diluya o se exprese de forma distinta en el conjunto. Cuarta, que la dosis efectiva de stevia difiera entre edades: los animales adultos consumieron, normalizado al peso corporal, menos de la mitad de la stevia que los juveniles (\S\ref{sec:res-consumo-tratamiento}), un patrón de exposición dependiente de la edad análogo al documentado para otros edulcorantes de bajas calorías \cite{tsan2022lcs}. Distinguir entre estas explicaciones excede el alcance de este trabajo, pero una replicación con mayor tamaño muestral y con la cohorte de origen identificada como factor explícito permitiría estimar cuánto de la inconsistencia observada es atribuible a cada una.

Ubicada en el eje inocuo-perjudicial que motiva el diseño de tres brazos (\S\ref{sec:bases}), la evidencia de esta tesis se inclina hacia el extremo no perjudicial: la stevia no reprodujo el fenotipo de la sacarosa, y los efectos aislados que se observaron en el NOR no se sostuvieron entre demoras ni entre edades. No obstante, la ausencia de un perjuicio consistente no equivale a una inocuidad demostrada: que la stevia sea un compuesto no calórico no garantiza que sea conductualmente inerte. En esa dirección apuntan tanto el hallazgo de que el consumo de edulcorantes de bajas calorías —incluida la stevia— durante la adolescencia se asocia a deterioros persistentes de la memoria hipocampal \cite{tsan2022lcs} como los trabajos previos del laboratorio anfitrión, que documentaron alteraciones persistentes de la fertilidad, la descendencia y la conducta tras la exposición temprana a stevia en hembras \cite{kruse2025stevia}; el efecto aislado de la stevia en T3 juvenil del NOR es, en ese sentido, puntualmente convergente con esa línea de evidencia, aunque su inconsistencia entre demoras impide extraer de él una conclusión firme. En el otro extremo, las propiedades antioxidantes y neuroprotectoras atribuidas a los glucósidos de esteviol \cite{pozdnyakova2025stevia,prata2017glycosides} tampoco permiten anticipar un resultado benigno sobre la memoria evaluada aquí: esos trabajos midieron peroxidación lipídica cerebral en ratas adultas con una dieta alta en grasas y sacarosa \cite{pozdnyakova2025stevia} o actividad insulino-mimética y antioxidante en fibroblastos cardíacos \textit{in vitro} \cite{prata2017glycosides}, desenlaces y modelos suficientemente distintos de la memoria de reconocimiento y la separación de patrones espaciales en ratas en desarrollo como para que una propiedad antioxidante periférica no se traduzca necesariamente en protección cognitiva. A esto se suma una limitación metodológica propia: la stevia se administró como un extracto acuoso de hojas (\S\ref{sec:soluciones}) que contiene, además de glucósidos de esteviol, fenoles, flavonoides y taninos (Tabla~\ref{tab:extracto}), y no como un compuesto purificado, lo que dificulta comparar directamente estos resultados con estudios que emplean compuestos purificados, como el de Prata et~al.\ \cite{prata2017glycosides}.

Otros estudios en roedores aportan matices en ambos sentidos. En modelos patológicos, el esteviósido tuvo un efecto antiamnésico en ratas tratadas con escopolamina \cite{singh2010stevioside}, la stevia modificó la plasticidad sináptica y los niveles de NADPH oxidasa en un modelo de desorden metabólico inducido por fructosa \cite{chavushyan2017stevia} y su administración revirtió la ansiedad y el deterioro de memoria en ratas diabéticas \cite{khakpai2023stevia}; en ratones, además, se comparó su efecto sobre la memoria y la histología hipocampal con el del aspartamo y la sucralosa \cite{villareal2016sweeteners}. Todos esos trabajos parten de contextos patológicos, de edulcorantes comparados o de otras especies y dosis, por lo que no son comparables de forma directa con el diseño de esta tesis, en el que animales sanos reciben la stevia en una ventana etaria definida y se evalúan tras un período de lavado. Tampoco puede descartarse que la stevia afecte los circuitos de recompensa de forma distinta que la sacarosa, como sugiere la menor expresión de $\Delta$FosB en el sistema dopaminérgico tras su consumo \cite{salinasvelarde2021fosb}; este aspecto no fue evaluado aquí.

\section{Conducta tipo ansiosa y exploración en el \textit{open field}}\label{sec:disc-of}

En el \textit{open field}, la distancia total recorrida no difirió entre grupos (\S\ref{sec:res-of-locomocion}), de modo que las diferencias en las demás variables no se explican, en principio, por un cambio de la locomoción general; esta lectura es menos firme para los grupos expuestos a stevia, que tuvieron 6 o 7 animales. Los animales expuestos a stevia presentaron más entradas, mayor distancia recorrida y más zonas visitadas en las zonas centrales y en la arena que los controles y que los expuestos a sacarosa, en ambas edades, un perfil que en esta prueba suele leerse como una menor conducta tipo ansiosa, dado que una menor exploración del centro indica un mayor nivel de ansiedad (\S\ref{sec:of}). La sacarosa, en cambio, se diferenció del control solo en la ventana juvenil, con menos entradas y menor distancia en las zonas 1+2, menos zonas visitadas y menos \textit{rearing}, un perfil compatible con una mayor conducta tipo ansiosa o con una menor exploración. En la dirección, estos resultados concuerdan con trabajos del laboratorio que describieron cambios duraderos de la conducta tipo ansiosa y de la memoria del miedo tras la exposición juvenil a sacarosa \cite{kruse2019sucrose} y alteraciones conductuales tras su consumo crónico \cite{rey2024metformin}, y contrastan con el hallazgo de que el consumo de sacarosa restringido a la adolescencia tardía no indujo conducta ansiosa pese a afectar la neurogénesis \cite{sanchezhuerta2022sucrose}, lo que sugiere, otra vez, que el efecto depende de la ventana de exposición.

Para la stevia, la literatura sobre conducta tipo ansiosa en animales sanos es escasa: un estudio en ratas diabéticas informó que su administración revirtió la ansiedad \cite{khakpai2023stevia}, un contexto metabólico distinto del de esta tesis, y los antecedentes del laboratorio en hembras describieron alteraciones conductuales de la descendencia tras la exposición temprana \cite{kruse2025stevia}. El perfil observado aquí no puede, por lo tanto, contrastarse con una referencia directa en animales sanos.

La interpretación de este perfil como una menor ansiedad requiere cautela por cuatro motivos. Primero, el \textit{open field} es una única prueba y no se complementó con otra de ansiedad (como el laberinto en cruz elevado, excluido de este diseño), de modo que corresponde hablar de conducta tipo ansiosa y no de un efecto ansiolítico. Segundo, el aumento de entradas y de zonas visitadas con stevia ocurrió sin cambio de la distancia total, lo que es compatible con más transiciones entre zonas y no solo con una menor aversión al centro. Tercero, el \textit{rearing}, una medida de exploración vertical, fue menor con stevia (en los juveniles frente al control y, en los adultos, frente a la sacarosa y, en el tiempo, frente al control), de modo que el perfil conductual no apunta en un único sentido: más exploración horizontal del centro y menos exploración vertical. Cuarto, los grupos stevia tuvieron 6 o 7 animales, el \textit{rearing} se puntuó manualmente por un único observador, y la interacción edad $\times$ tratamiento fue significativa para algunas variables (distancia en las zonas 1+2, zonas visitadas y número de eventos de \textit{rearing}) y no para otras (entradas a las zonas 1+2, distancia en la zona 1 y tiempo de \textit{rearing}), por lo que no puede afirmarse una dependencia de la edad general. Estos resultados se consideran, por todo ello, exploratorios.

El perfil del \textit{open field} agrega una dimensión al contraste con la hipótesis (\S\ref{sec:disc-sintesis}): la stevia no reprodujo en esta prueba el perfil asociado a la sacarosa juvenil, lo que concuerda con la primera mitad de la hipótesis, mientras que su posición respecto del control ---con más exploración central pero menos \textit{rearing}--- no permite concluir una condición equivalente o superior.

\section{Ventana etaria}\label{sec:disc-edad}

El papel de la ventana etaria de exposición difiere notablemente entre las dos pruebas. En el NOR, la interacción edad $\times$ tratamiento fue significativa en las tres demoras (Tabla~\ref{tab:nor-anova}), lo que permite afirmar que el efecto de al menos uno de los edulcorantes dependió de la edad de exposición. En la SLR, en cambio, esa misma interacción no fue significativa en ninguna configuración (Tabla~\ref{tab:slr-anova}), por lo que, pese a que los valores puntuales de DI sugieren un patrón distinto entre juveniles y adultos (\S\ref{sec:disc-sacarosa}), los datos no sostienen estadísticamente una dependencia de la edad para la separación de patrones espaciales en este diseño.

La interpretación del factor edad está, además, condicionada por un rasgo del diseño: "edad" confunde aquí la ventana de exposición con la edad del animal al momento del testeo, ya que la batería conductual se administró recién al finalizar el lavado, es decir, alrededor de PD75 en el grupo juvenil y de PD125 en el adulto (\S\ref{sec:diseno}). Un efecto atribuido a mayor vulnerabilidad juvenil podría reflejar, en principio, tanto la inmadurez del circuito en el momento de la exposición como alguna diferencia asociada a la edad en el momento de la prueba; este diseño no permite separar ambos componentes, y la discusión de vulnerabilidad diferencial debe leerse con esa salvedad.

Por último, el índice de preferencia calculado en \S\ref{sec:res-consumo-tratamiento} es relevante para interpretar la dosis efectiva por edad: pese a que los animales juveniles consumieron considerablemente más solución endulzada que los adultos cuando ese consumo se normaliza al peso corporal, la preferencia por la solución endulzada frente al agua fue similar entre edades, y la misma diferencia de magnitud por edad se observó también en el consumo de agua de los controles. Esto sugiere que la mayor ingesta normalizada de los juveniles refleja una diferencia de escala corporal propia de animales en crecimiento más que una mayor motivación por el sabor dulce, y respalda tratar la diferencia de dosis efectiva entre edades —invocada en \S\ref{sec:disc-stevia} como una de las explicaciones posibles de la inconsistencia de la stevia en el NOR— como un factor de exposición y no de preferencia.

\section{Limitaciones}\label{sec:disc-limitaciones}

Varias limitaciones del diseño condicionan el alcance de estas conclusiones. Los grupos control de la SLR provienen exclusivamente de las cohortes corridas en 2025, mientras que los grupos stevia y sacarosa de 2026, evaluados solo en la SLR, no tuvieron un control contemporáneo (\ref{anx:muestra}); si existiera una deriva temporal en el manejo de los animales, el equipamiento o el observador entre años, esta podría tanto generar diferencias espurias entre tratamiento y control como enmascarar diferencias reales, en cualquiera de los dos sentidos, lo que añade incertidumbre a todas las comparaciones de la SLR que incluyen cohortes de 2026 y no solo a las de stevia. El tamaño muestral de 9 a 10 animales por grupo (\S\ref{sec:res-nor}, \S\ref{sec:res-slr}), estimado a partir de estudios previos del laboratorio \cite{charan2013sample}, ofrece potencia razonable para los efectos de sacarosa, de magnitud comparable a la de esos estudios previos, pero probablemente resulta insuficiente para detectar efectos pequeños o intermedios de la stevia; esto sesga el diseño hacia no detectar un efecto real de la stevia (error de tipo~II) antes que hacia concluir con confianza su ausencia, una distinción central dado que buena parte de la discusión de la stevia descansa en resultados no significativos. En sentido inverso, el análisis de cada demora del NOR y cada configuración de la SLR mediante ANOVA independientes, sin una corrección conjunta para el conjunto de la batería, incrementa el riesgo de que alguno de los hallazgos puntuales de la stevia —en particular T2 adulto y T3 juvenil del NOR— sea un falso positivo, en la dirección opuesta a la limitación anterior.

A esto se suman limitaciones de carácter estadístico y de diseño más acotadas. Los residuos del modelo de s-SLR se apartaron de la normalidad, lo que llevó a confirmar las diferencias significativas mediante Mann-Whitney (\S\ref{sec:disc-sacarosa}), pero no se dispone de una alternativa no paramétrica para el efecto principal y la interacción del ANOVA de esa configuración. No se evaluó un efecto de cohorte antes de combinar los datos de las distintas cohortes (\S\ref{sec:diseno}), de modo que una diferencia entre cohortes no puede separarse del efecto de los tratamientos. El estudio se realizó exclusivamente en machos, de modo que ninguna de estas conclusiones es extrapolable a hembras, población en la que el mismo laboratorio documentó efectos de la stevia sobre la fertilidad, la descendencia y la conducta \cite{kruse2025stevia}. La exploración del día 1 de la SLR, además, no fue equivalente entre grupos (\S\ref{sec:res-slr-exploracion}), por lo que no puede descartarse que la codificación difiriera entre ellos. En el \textit{open field}, los grupos expuestos a stevia tuvieron 6 o 7 animales, el número de animales difiere entre variables y el \textit{rearing} se puntuó manualmente por un único observador, sin estimación de concordancia entre observadores (\S\ref{sec:res-of}). Finalmente, el análisis histológico y el control de exploración en T1 del NOR todavía no están disponibles; mientras estos sigan pendientes, ninguna interpretación mecanística de los hallazgos conductuales puede considerarse más que una hipótesis a evaluar.

\section{Integración pendiente}\label{sec:disc-integracion}

La resolución de buena parte de las ambigüedades señaladas en este capítulo —en particular, si los grupos stevia constituyen un perfil homogéneo o si su variabilidad refleja una mezcla de animales con trayectorias distintas— depende de análisis que todavía no se han incorporado a esta tesis. La cuantificación de PCNA y DCX (\S\ref{sec:res-neurogenesis}) podría indicar en qué etapa de la cascada neurogénica, si en alguna, inciden la sacarosa y la stevia, y si el perfil histológico de la stevia es intermedio o heterogéneo de un modo consistente con su perfil conductual. El análisis multivariado y de aprendizaje automático (\S\ref{sec:res-ml}, objetivos M1 y M2) podría, por su parte, indicar si los animales expuestos a stevia emergen como un grupo conductual diferenciado o se distribuyen entre los perfiles de control y de sacarosa, una pregunta que los contrastes univariados de este capítulo no están en condiciones de responder. El desarrollo de estas líneas se retoma en el Capítulo~\ref{cap:perspectivas}; aquí quedan señaladas como el paso necesario para pasar de un resultado no concluyente a una conclusión sobre la stevia.

\pendiente{Esta Discusión cubre NOR, SLR y \textit{open field}. Falta discutir: consumo y peso corporal (\S\ref{sec:res-consumo}), neurogénesis/histología (\S\ref{sec:res-neurogenesis}) y análisis multivariado/aprendizaje automático (\S\ref{sec:res-ml}), una vez completados esos resultados.}
